%% kapa: exam=1 round=30 year=2019 period=1 form=A subject=economics subject_ko=경제학원론 seq=17 number=57
%% source: Q-Net 감정평가사 기출문제 (cst00309) artlSeq=5205205 file=extracted/1008/5205205/2019년 제30회 감정평가사 1차시험 문제지/2019년 제30회 감정평가사 1차시험 문제지/2019년도 제30회 감정평가사 1차 1교시 A형.hwp
%% conversion: kordoc 4.24.51 -> tools/gen1.py (auto); review=proofread-sonnet; flags=none
\begin{question}{57}
정상재 A, B의 가격이 각각 2 \% 상승할 때 A재의 소비지출액은 변화가 없었지만, B재의 소비지출액은 1 \% 감소하였다. 이 때 두 재화에 대한 수요의 가격탄력성 ${ \epsilon_{ A } }$, ${ \epsilon_{ B } }$에 관한 설명으로 옳은 것은? (단, ${ \epsilon_{ A } }$와 ${ \epsilon_{ B } }$는 절댓값으로 표시한다.)
\begin{choices}
  \item $\epsilon_{ A } >1, \epsilon_{ B } >1$
  \item $\epsilon_{ A } =1, \epsilon_{ B } >1$
  \item $\epsilon_{ A } =0, \epsilon_{ B } <1$
  \item $\epsilon_{ A } =1, \epsilon_{ B } <1$
  \item $\epsilon_{ A } <1, \epsilon_{ B } <1$
\end{choices}
\end{question}
\answer{2}
